% Generated draft: requires mathematical and source review.
\documentclass[11pt]{article}
\usepackage[margin=25mm]{geometry}
\usepackage{fontspec}
\setmainfont{Latin Modern Roman}
\usepackage{amsmath,amssymb,mathtools,unicode-math}
\setmathfont{Latin Modern Math}
\usepackage{enumitem,needspace,xurl,hyperref}
\setlength{\parindent}{0pt}
\setlength{\parskip}{5pt}
\setlength{\emergencystretch}{3em}
\newcommand{\N}{\mathbb{N}}
\newcommand{\Z}{\mathbb{Z}}
\newcommand{\Q}{\mathbb{Q}}
\newcommand{\R}{\mathbb{R}}
\newcommand{\C}{\mathbb{C}}
\newcommand{\F}{\mathbb{F}}
\newcommand{\A}{\mathbb{A}}
\newcommand{\PP}{\mathbb{P}}
\newcommand{\E}{\mathbb{E}}
\newcommand{\eps}{\varepsilon}
\newcommand{\dd}{\,\mathrm{d}}
\begin{document}
\section*{UnsolvedMath ID: 30004852}
\textbf{Source problem code:} catalogue entry
\subsection{Definitions and mathematical statement}
For integers \(n\geq 1\) and \(d\geq 1\), and a polynomial or holomorphic germ \(f\colon (\C^n,0)\to(\C,0)\), the displayed jet-space definition in the supplied record can be reconstructed as
\[
\chi_d(f)=\left\{u(t)=\sum_{i=1}^{d}a_i t^i\;:\;a_i\in\C^n,\quad f(u(t))\equiv t^d\pmod{t^{d+1}}\right\}.
\]
Thus \(\chi_d(f)\) consists of jets centered at the origin whose composite with \(f\) has order exactly \(d\) and leading coefficient \(1\). For a germ, composition is interpreted using its Taylor expansion modulo \(t^{d+1}\). Write \(H_c^j\) for compactly supported cohomology in degree \(j\); the coefficient ring is not specified in the record.

The source attributes to \([\mathrm{BdBLN20}]\) an arc--Floer conjecture whose displayed formula is recoverable as
\[
HF^k(\phi_f^{\,d})\cong H_c^{\,k+2nd+n-1}(\chi_d(f)),\qquad k\in\Z.
\]
Here \(\phi_f^{\,d}\) denotes the \(d\)-fold iterate of the map denoted by \(\phi_f\) in the source. The intended claim concerns every positive iterate \(d\) and every \(f\) satisfying the original setup. That setup is absent from the supplied excerpt: neither \(\phi_f\), the Floer theory \(HF^*\), its grading and coefficient conventions, nor the complete hypotheses on \(f\) are defined. In particular, the excerpt does not establish whether the polynomial and isolated-singularity hypotheses stated later also govern this conjecture. The displayed isomorphism is therefore a conditional reconstruction of the intended statement, not a fully specified conjecture. Its degree shift is transcribed from the supplied text and requires comparison with the original typeset formulation.

The report says that the authors of \([\mathrm{BdBLN20}]\) verified this conjecture when \(d\) equals the multiplicity of \(f\) at \(0\). Immediately afterward, it labels the general assertion ``Theorem 3 (M., in progress).'' Neither this announcement nor the supplied verification note is treated here as a verified proof or a current literature determination; the identity represented by ``M.'' is not supplied.

The proposed proof construction asks for a natural PSS morphism
\[
\operatorname{ev}\colon HF^k(\phi_f^{\,d})\longrightarrow H_c^{\,k+2nd+n-1}(\chi_d(f)).
\]
At chain level, the source proposes using a moduli space of holomorphic disks whose boundaries limit to an orbit in the mapping torus of \(\phi_f^{\,d}\), with their \(d\)-jets sweeping out a cycle in a thickening of \(\chi_d(f)\) inside the space of centered \(d\)-jets, denoted there by \(\operatorname{Jet}_d(f)|_0\). It then asks for filtrations \(F\) and \(F'\) on the respective complexes such that this map preserves the filtrations and induces an isomorphism on associated graded objects. The source associates these filtrations with spectral sequences in \([\mathrm{McL19}]\) and \([\mathrm{BdBLN20}]\), respectively, and compares the proposed map with the logarithmic PSS map of Ganatra and Pomerleano \([\mathrm{GP20}]\). The analytic data, precise jet-space notation, filtrations, and spectral-sequence convergence hypotheses are not supplied.

The excerpt also poses further questions. It describes a pair-of-pants product on the bigraded group
\[
\bigoplus_{d\geq1}HF^*(\phi_f^{\,d})
\]
and asks for its counterpart on \(\bigoplus_{d\geq1}H_c^*(\chi_d(f))\). The iterate range is interpreted here as the positive integers used for the jet loci; the source writes only \(\bigoplus_d\). It suggests a possible relation to an arc-space version of the string product, with arc boundaries considered as loops, and possible structure involving a framed little disks operad and other operations. These are suggestions, not specified algebraic structures. In particular, the product on the cohomology side and the grading adjustments induced by the \(d\)-dependent shift remain to be defined.

Finally, for polynomials \(f_0,f_1\) with isolated singularities at the origin, the source proposes functoriality under adjacency: if \(f_1\) is a sufficiently small perturbation of \(f_0\), there should be a corresponding algebra morphism
\[
\bigoplus_{d\geq1}H_c^*(\chi_d(f_0))\longrightarrow
\bigoplus_{d\geq1}H_c^*(\chi_d(f_1)).
\]
The phrase ``sufficiently small perturbation,'' the deformation data, the relevant topology, the products, and the precise functoriality requirements are unspecified. The supplied record does not identify which of the isomorphism, its proposed construction, the product question, or adjacency functoriality is intended as the standalone open target.
\subsection{Short English statement}
The record concerns a proposed degree-shifted identification between Floer cohomology of iterated maps associated with a singularity and compactly supported cohomology of normalized jet contact loci. It also asks about products, additional operations, and maps under adjacency. Missing setup and conflicting supplied status descriptions prevent a fully specified open-problem statement; this is a statement-only reconstruction requiring human mathematical review.
\subsection{Sources}
\begin{itemize}
\item Local UnsolvedMath registry, permanent ID \(30004852\), title ``Floer Cohomology and Contact Loci of Jets,'' exact source code \texttt{catalogue entry}; supplied complete source record. The project identifies \texttt{data/unsolvedmath/problems.json} and \texttt{catalog/problems.json} as registry collections; the record's location in either collection has not been independently checked. No proposer is named; the supplied proposed year is \(2021\).
\item Supplied report citation: \emph{Singularities (2022)}, Oberwolfach Reports, pp.\ 2573--2639, DOI \texttt{10.4171/owr/2021/47}. Source URL and literature-source link labeled ``original OWR formulation'': \url{https://doi.org/10.4171/owr/2021/47}. The record describes Oberwolfach Reports as documenting mathematical workshops and their open-problem discussions.
\item Supplied literature-source link: Budur--Fern\'andez de Bobadilla--L\^e--Nguyen, ``cohomology of contact loci'': \url{https://doi.org/10.4310/jdg/1649953456}. The excerpt attributes Conjecture 2 and the multiplicity case to \([\mathrm{BdBLN20}]\); the exact bibliographic correspondence requires verification.
\item Supplied literature-source link: de Lorenzo Poza--Huang, ``semihomogeneous contact loci'': \url{https://arxiv.org/abs/2506.18805}.
\item Additional references appearing only as abbreviated labels in the excerpt: Ganatra and Pomerleano \([\mathrm{GP20}]\), for a logarithmic PSS map; \([\mathrm{McL19}]\), for a Floer spectral sequence; \([\mathrm{Sei08}]\), Section 2, for the pair-of-pants product; \([\mathrm{CS04}]\), for the string product; \([\mathrm{Vor05}]\), for the framed little disks operad; and \([\mathrm{Sei19}]\), for other structures. Complete bibliographic records are not supplied. ``M.'' is the only supplied attribution for the announced theorem in progress.
\item Literature assessment dated exactly \texttt{2026-08-22}: this is supplied dated curation, not independently verified. It labels the problem partially solved and says that the general arc--Floer isomorphism remains open, while the multiplicity case, important special families, hyperplane arrangements, and 2025 semihomogeneous results provide further cases. The background explicitly calls this a dated literature triage derived from the supplied OWR curation and requests independent verification of time-sensitive status and cited claims.
\item The supplied statement-verification note says that the official contribution announces the general isomorphism as a theorem in progress immediately after the conjecture, so that contribution does not present it there as a still-open problem. The record marks the statement \texttt{unrecoverable}, supplies no clean statement, and retains the original record for provenance. These are supplied editorial assessments, not independently verified conclusions.
\end{itemize}
\end{document}
